\documentclass[11pt,a4paper]{article}

% === 中文支持 ===
\usepackage{ctex}
\usepackage[T1]{fontenc}

% === 页面布局 ===
\usepackage[top=2.5cm, bottom=2.5cm, left=2.5cm, right=2.5cm]{geometry}

% === 表格 ===
\usepackage{booktabs}
\usepackage{array}
\usepackage{longtable}
\usepackage{enumitem}

% === 代码 ===
\usepackage{listings}
\lstset{basicstyle=\ttfamily\footnotesize, breaklines=true, frame=single, showstringspaces=false}

% === 颜色 ===
\usepackage{xcolor}
\definecolor{primary}{HTML}{404040}
\definecolor{success}{HTML}{22c55e}
\definecolor{danger}{HTML}{e74c3c}
\definecolor{info}{HTML}{3b82f6}

% === 超链接 ===
\usepackage[hidelinks]{hyperref}
\hypersetup{colorlinks=true, linkcolor=primary, urlcolor=info}

% === 标题样式 ===
\usepackage{titlesec}
\titleformat{\section}{\Large\bfseries\color{primary}}{}{0em}{}[\titlerule]
\titleformat{\subsection}{\large\bfseries\color{primary}}{}{0em}{}
\titleformat{\subsubsection}{\normalsize\bfseries\color{primary}}{}{0em}{}

% === 文档信息 ===
\title{\textbf{项目文档规范}\\[0.5em]\small 文档命名、格式与生命周期管理}
\author{万能产品小助手（PM AI）}
\date{2026-05-06 / 版本 v0.1}

\begin{document}

\maketitle

\tableofcontents
\newpage

% ============================================================
\section{总则}
% ============================================================

本项目所有正式文档采用 \textbf{tex-pdf} 格式管理，确保结构规范性、版本可追溯性和长期可读性。

\begin{itemize}
  \item \textbf{源文件}：\texttt{.tex} 格式，可 Git 版本控制
  \item \textbf{交付物}：\texttt{.pdf} 格式，用于审阅和分发
  \item \textbf{编译工具}：Tectonic（\texttt{\textasciitilde/.local/tectonic/tectonic}）
  \item \textbf{文档管理者}：PM AI（万能产品小助手），负责文档的创建、更新、编译和归档
\end{itemize}

\subsection{版本阶段定义}

\begin{table}[h]
\centering
\begin{tabular}{clp{7cm}}
\toprule
\textbf{阶段} & \textbf{名称} & \textbf{内容范围} \\
\midrule
V0 & 先备阶段 & 智能体分工、商店整体框架、文档体系搭建、技术选型确认 \\
V0.x & 先备迭代 & 分工调整、架构调研、参考方案分析、规范补充 \\
V1 & 正式需求 & 第一批正式业务需求接入 \\
V1.x & 功能迭代 & 功能增强、Bug 修复、性能优化 \\
V2+ & 演进阶段 & 架构升级、生态扩展、商业运营 \\
\bottomrule
\end{tabular}
\caption{版本阶段定义}
\end{table}

\textbf{当前阶段}：V0（先备阶段）

% ============================================================
\section{目录结构}
% ============================================================

\begin{lstlisting}[basicstyle=\ttfamily\footnotesize, frame=single]
enterprise-app-store/
  core/                               # 核心产品文档（MRD/PRD/RFC/TRD）
    MRD.tex/pdf                     # 市场需求文档
    PRD.tex/pdf                     # 产品需求文档
    RFC.tex/pdf                     # 请求意见/技术提案
    TRD.tex/pdf                     # 技术需求文档
  reference/                          # 外部调研文档（仅放外部参考对象的架构调研）
    DOC-NAMING-CONVENTION.tex/pdf    # 本文件：文档命名规范
    REF-001-skill-store-arch-survey.tex/pdf
    REF-002-*.tex/pdf                # 项目总纲
    REF-003-open-skills-ecosystem.tex/pdf # skills.sh 开放技能生态架构调研
    ...
  demands/                            # 需求与项目规范文档
    MASTER.tex/pdf                    # 总需求文档（唯一，持续更新）
    DEMAND-000-v0-master.tex/pdf      # V0 阶段需求快照（归档）
    DEMAND-001-role-redefine.tex/pdf  # 增量需求文档
    DEMAND-002-dev-env-setup.tex/pdf # 开发环境搭建
    DEMAND-003-task-workflow.tex/pdf # 任务流转路径
    DEMAND-004-multi-agent-collaboration.tex/pdf  # 多智能体协作
    GUIDE-001-project-structure.tex/pdf  # 项目目录结构与开发规范
    GUIDE-002-project-methodology.tex/pdf  # 项目方法论总纲（智能体搭配+选型+文档+自闭环）
    ...
\end{lstlisting}

\subsection{目录职责}

\begin{table}[h]
\centering
\begin{tabular}{lp{9cm}}
\toprule
\textbf{目录} & \textbf{职责} \\
\midrule
\texttt{core/} & 核心产品文档。MRD（市场需求）、PRD（产品需求）、RFC（技术提案）、TRD（技术需求）。这是产品从市场分析到技术落地的完整文档链，PM AI 负责维护。 \\
\texttt{reference/} & 仅存放外部调研文档。包括外部项目架构调研、竞品分析、开源方案研究等。每个外部参考对象独立编号（REF-001, REF-002, ...）。注意：项目自身的规范、指南、总纲等不属于此目录。 \\
\texttt{demands/} & 存放需求文档和项目规范文档。包含三类：
\textbf{(1) 总需求文档}（MASTER.tex/pdf）：项目全部需求的汇总，持续更新。
\textbf{(2) 增量需求文档}（DEMAND-xxx.tex/pdf）：每次需求新增或变更的独立文档。
\textbf{(3) 项目规范文档}（GUIDE-xxx.tex/pdf）：目录结构、编码规范、开发流程等项目内部规范。 \\
\bottomrule
\end{tabular}
\end{table}

% ============================================================
\section{文件命名规范}
% ============================================================

\subsection{命名格式}

所有文档文件采用以下命名格式：

\begin{quote}
\texttt{\{类型前缀\}-\{三位编号\}-\{kebab-case 简述\}.\{tex|pdf\}}
\end{quote}

\subsection{需求文档的两层结构}

需求文档采用	extbf{总需求文档 + 增量需求文档}两层结构：

\begin{enumerate}
  \item \textbf{总需求文档}（\texttt{MASTER.tex/pdf}）：
    \begin{itemize}[nosep]
      \item 项目唯一的、持续更新的需求汇总
      \item 每次新增/变更需求后，同步更新 MASTER
      \item 内容包含：需求总表（编号、名称、优先级、状态）、每个需求的摘要
      \item 版本号随每次更新递增
      \item 类似于「活的」需求 Backlog
    \end{itemize}
  \item \textbf{增量需求文档}（\texttt{DEMAND-xxx.tex/pdf}）：
    \begin{itemize}[nosep]
      \item 每次需求新增或变更时生成的独立文档
      \item 记录该次需求的完整描述、决策、交付物、验收结论
      \item 编号递增（DEMAND-001, DEMAND-002, ...）
      \item 一旦归档，不再修改（变更通过新 DEMAND 记录）
      \item 类似于「冻结的」变更记录
    \end{itemize}
\end{enumerate}

\textbf{两者的关系}：DEMAND-xxx 是每次需求的详细快照，MASTER 是所有需求的汇总视图。每次生成 DEMAND-xxx 后，必须同步更新 MASTER。

\subsection{类型前缀}

\begin{table}[h]
\centering
\begin{tabular}{lll}
\toprule
\textbf{前缀} & \textbf{含义} & \textbf{存放目录} \\
\midrule
\texttt{REF} & 外部调研文档（Reference） & \texttt{reference/} \\
\texttt{DEMAND} & 增量需求文档（Demand） & \texttt{demands/} \\
\texttt{MASTER} & 总需求文档 & \texttt{demands/} \\
\texttt{GUIDE} & 项目规范文档（Guide） & \texttt{demands/} \\
\texttt{SPEC} & 技术规格文档（Specification） & \texttt{demands/} \\
\bottomrule
\end{tabular}
\end{table}

\subsection{命名示例}

\begin{table}[h]
\centering
\begin{tabular}{ll}
\toprule
\textbf{文件名} & \textbf{说明} \\
\midrule
\texttt{REF-001-skill-store-arch-survey} & 第 1 份参考文档：技能商店架构调研 \\
\texttt{REF-002-tex-pdf-best-practices} & 第 2 份参考文档：tex-pdf 最佳实践 \\
\texttt{DEMAND-001-role-redefine} & 第 1 份需求：角色分工重新定义 \\
\texttt{DEMAND-002-search-engine-integration} & 第 2 份需求：搜索引擎集成 \\
\texttt{SPEC-001-api-gateway-design} & 第 1 份技术规格：API 网关设计 \\
\texttt{GUIDE-001-deployment-guide} & 第 1 份指南：部署操作指南 \\
\bottomrule
\end{tabular}
\end{table}

\subsection{命名规则}

\begin{enumerate}
  \item \textbf{编号全局递增}：REF 和 DEMAND 各自独立编号，不交叉
  \item \textbf{简述用 kebab-case}：单词间用连字符连接，全小写，无空格
  \item \textbf{简述控制在 2--5 个单词}：够区分即可，不要太长
  \item \textbf{同一需求多次更新}：不修改文件名，更新 tex 源文件内容后重新编译 pdf
  \item \textbf{tex 和 pdf 同名}：\texttt{DEMAND-001-role-redefine.tex} 对应 \texttt{DEMAND-001-role-redefine.pdf}
\end{enumerate}

% ============================================================
\section{tex-pdf 模板规范}
% ============================================================

\subsection{统一文档头}

每份 tex 文档必须包含以下文档头：

\begin{lstlisting}[basicstyle=\ttfamily\footnotesize, frame=single]
\documentclass[11pt,a4paper]{article}
\usepackage{ctex}                        % 中文支持（必须）
\usepackage[T1]{fontenc}
\usepackage[top=2.5cm, bottom=2.5cm,
            left=2.5cm, right=2.5cm]{geometry}
\usepackage{booktabs}                    % 表格（推荐）
\usepackage{amssymb}                     % 数学符号
\usepackage{hyperref}                    % 超链接
\usepackage{titlesec}                    % 标题样式

\title{\textbf{标题}}
\author{万能产品小助手（PM AI）\\魏子政 审阅}
\date{YYYY-MM-DD / 版本 vX.Y}
\end{lstlisting}

\subsection{通用编译命令}

\begin{lstlisting}[basicstyle=\ttfamily\footnotesize, frame=single]
# 编译（自动下载依赖、自动多次运行处理交叉引用）
cd <对应目录> && ~/.local/tectonic/tectonic <文件名>.tex

# 示例
cd reference && ~/.local/tectonic/tectonic REF-001-xxx.tex
cd demands && ~/.local/tectonic/tectonic DEMAND-001-xxx.tex
\end{lstlisting}

\subsection{已知限制与避坑指南}

\begin{enumerate}
  \item \textbf{避免 fontawesome5 命令}：\texttt{\textbackslash faCalendar} 等命令不可用，用文字或 amssymb 符号替代
  \item \textbf{避免 Unicode 箭头}：\texttt{->} 等特殊字符需要用 \texttt{\textbackslash to} 或 \texttt{-{}->{}} 替代
  \item \textbf{避免 listings 语言高亮}：Tectonic 可能缺少部分语言定义，使用无 language 参数的 lstlisting
  \item \textbf{颜色定义}：直接用 \texttt{\textbackslash definecolor}，不要依赖 xcolor 的命名色
  \item \textbf{Overfull hbox}：少量警告可接受，不影响阅读
\end{enumerate}

% ============================================================
\section{文档生命周期}
% ============================================================

\subsection{创建流程}

\begin{enumerate}
  \item PM AI 识别需要创建文档的场景（需求变更、外部调研、技术预研等）
  \item 确定文档类型（REF / DEMAND / SPEC / GUIDE）和编号
  \item 按命名规范创建 tex 文件
  \item 编写内容
  \item 编译为 pdf
  \item 通过飞书发送给魏子政审阅
\end{enumerate}

\subsection{更新流程}

\begin{enumerate}
  \item 需求变更或内容更新触发
  \item 更新 tex 源文件中的日期和版本号
  \item 在文档末尾追加更新记录
  \item 重新编译 pdf
  \item 通知魏子政审阅
\end{enumerate}

\subsection{更新记录格式}

每份文档末尾维护更新记录：

\begin{lstlisting}[basicstyle=\ttfamily\footnotesize, frame=single]
\section*{更新记录}
\begin{tabular}{llll}
\toprule
版本 & 日期 & 变更内容 & 作者 \\
\midrule
v0.1 & 2026-05-06 & 初始版本 & PM AI \\
v0.2 & 2026-05-07 & 新增第三节 & PM AI \\
\bottomrule
\end{tabular}
\end{lstlisting}

% ============================================================
\section{目录管理方式}
% ============================================================

本文档由 PM AI（万能产品小助手）负责维护。以下是对 \texttt{demands/} 和 \texttt{reference/} 两个目录的管理规则，确保后续迭代有据可依。

\subsection{reference/ --- 外部调研文档管理}

\subsubsection{目录定位}

\texttt{reference/} 仅存放\textbf{外部调研}产生的文档。判断标准：

\begin{itemize}[nosep]
  \item \textbf{放入 reference/}：外部项目架构分析、竞品调研、开源方案研究、行业标准参考
  \item \textbf{不放入 reference/}：项目自身的规范、指南、总纲、需求文档（这些放 \texttt{demands/} 或 \texttt{core/}）
  \item \textbf{不放入 reference/}：技术选型文档、团队协作文档（这些放项目根目录 \texttt{.md}）
\end{itemize}

\subsubsection{编号规则}

\begin{itemize}[nosep]
  \item 前缀固定为 \texttt{REF}
  \item 编号三位递增：\texttt{REF-001}、\texttt{REF-002}、\texttt{REF-003}、...
  \item 编号\textbf{全局唯一}，不重复使用
  \item 文件名格式：\texttt{REF-\{NNN\}-\{kebab-case\}.tex/pdf}
\end{itemize}

\subsubsection{创建流程}

\begin{enumerate}[nosep]
  \item PM AI 收到外部调研任务（魏子政指派或主动调研）
  \item 确定下一个可用编号（查看本文件索引表中最后一个 REF 编号 +1）
  \item 创建 \texttt{REF-\{NNN\}-\{kebab-case\}.tex}，使用统一 tex 模板
  \item 编写调研内容
  \item 编译为 pdf：\texttt{cd reference \&\& \textasciitilde/.local/tectonic/tectonic REF-\{NNN\}-\{kebab-case\}.tex}
  \item 更新本文件的索引表（添加新行）
  \item 重新编译本文件为 pdf
  \item Git commit
\end{enumerate}

\subsubsection{更新规则}

\begin{itemize}[nosep]
  \item REF 文档\textbf{创建后一般不修改}（外部调研结果是快照）
  \item 如需补充信息：更新 tex 源文件 → 递增版本号 → 追加更新记录 → 重新编译
  \item 如有新调研对象：创建新的 REF-xxx（不修改已有 REF）
\end{itemize}

\subsubsection{与 DEMAND 的关系}

REF 文档是\textbf{输入}，DEMAND 文档是\textbf{输出}。典型链路：
\begin{verbatim}
外部调研 → REF-xxx（记录调研结果）
         → DEMAND-xxx（基于调研结论提出项目需求）
         → MASTER（更新需求总表）
\end{verbatim}

\subsection{demands/ --- 需求与规范文档管理}

\subsubsection{目录定位}

\texttt{demands/} 存放三类文档：

\begin{enumerate}[nosep]
  \item \textbf{MASTER}：项目总需求文档（唯一，持续更新）
  \item \textbf{DEMAND-xxx}：增量需求文档（每个需求独立编号，归档后不修改）
  \item \textbf{GUIDE-xxx}：项目规范文档（编码规范、流程指南等）
\end{enumerate}

\subsubsection{编号规则}

\begin{itemize}[nosep]
  \item MASTER 无编号（唯一文件）
  \item DEMAND 和 GUIDE 各自独立编号，三位递增
  \item 编号全局唯一，不重复使用
  \item 文件名格式：\texttt{\{PREFIX\}-\{NNN\}-\{kebab-case\}.tex/pdf}
\end{itemize}

\subsubsection{MASTER 的管理方式}

MASTER 是项目的\textbf{唯一需求权威来源}：

\begin{itemize}[nosep]
  \item 每次新增需求后\textbf{必须}更新 MASTER（添加新行到需求总表）
  \item 每次需求状态变更后\textbf{必须}更新 MASTER（修改状态字段）
  \item 版本号随每次更新递增（v1.0 → v1.1 → v1.2 → ...）
  \item MASTER 中的需求编号（D-001, D-002, ...）与 DEMAND 文件编号对应但不完全相同（一个 D-xxx 可能拆为多个 DEMAND-xxx，也可能合并）
  \item 更新后必须重新编译 pdf
\end{itemize}

\subsubsection{DEMAND-xxx 的管理方式}

DEMAND 是需求的\textbf{冻结快照}：

\begin{itemize}[nosep]
  \item 创建时：确定编号 → 编写完整需求 → 编译 pdf → 更新 MASTER → Git commit
  \item 归档后：\textbf{不再修改}原文档
  \item 需求变更：创建新的 DEMAND-xxx（引用旧编号，说明变更内容），不修改旧文档
  \item 每份 DEMAND 包含：需求描述、拍板决策、技术方案、验收标准、状态
\end{itemize}

\subsubsection{GUIDE-xxx 的管理方式}

GUIDE 是项目的\textbf{持续维护规范}：

\begin{itemize}[nosep]
  \item GUIDE 与 DEMAND 不同，GUIDE 可以持续更新（规范会随项目演进）
  \item 更新时递增版本号，追加更新记录
  \item GUIDE-002（项目方法论总纲）是核心文档，任何方法论变更都必须同步更新
\end{itemize}

\subsubsection{创建流程（DEMAND 和 GUIDE 通用）}

\begin{enumerate}[nosep]
  \item PM AI 识别需求/规范创建场景
  \item 确定类型（DEMAND 或 GUIDE）和下一个可用编号
  \item 创建 tex 文件，使用统一模板
  \item 编写内容（DEMAND 需包含验收标准，GUIDE 需包含操作指引）
  \item 编译为 pdf：\texttt{cd demands \&\& \textasciitilde/.local/tectonic/tectonic \{文件名\}.tex}
  \item 更新 MASTER（DEMAND 必须更新，GUIDE 视情况）
  \item 更新本文件的索引表
  \item 重新编译本文件为 pdf
  \item Git commit
\end{enumerate}

\subsection{版本迭代与阶段划分}

\subsubsection{V0 先备阶段}

V0 阶段的目标是建立项目基础设施：

\begin{itemize}[nosep]
  \item 智能体分工定义（OpenClaw + Cursor）
  \item 技术选型与锁定
  \item 文档体系建立
  \item 开发环境搭建
  \item MVP 核心流程跑通
\end{itemize}

V0 阶段对应的文档：DEMAND-000（V0 总纲）、DEMAND-001 至 DEMAND-004（分工/环境/流程/协作）、GUIDE-001（结构）、GUIDE-002（方法论）。

\subsubsection{V1 正式需求阶段}

V1 起进入正式需求迭代：

\begin{itemize}[nosep]
  \item DEMAND-005 起（MVP 功能）开始编号
  \item 每个正式需求走完整流程：记录 → 评估 → 设计 → 拆解 → 开发 → 验收 → 文档更新
  \item MASTER 持续更新，记录所有需求的最新状态
  \item REF 文档按需新增（新调研对象）
\end{itemize}

\subsubsection{跨阶段文档处理}

\begin{itemize}[nosep]
  \item V0 文档归档保留，不删除（历史追溯需要）
  \item V0 的方法论文档（GUIDE-002）持续维护，不随阶段归档
  \item MASTER 跨阶段持续更新，V1 的需求追加到 V0 需求之后
  \item core/ 文档（MRD/PRD/RFC/TRD）跨阶段持续更新
\end{itemize}

\subsection{索引维护规则}

本文件（DOC-NAMING-CONVENTION.tex）是\textbf{文档索引的唯一权威来源}：

\begin{itemize}[nosep]
  \item 每次新增/归档/重命名文档后，必须更新对应索引表
  \item 索引表中的状态字段必须与实际一致
  \item 更新后重新编译 pdf
\end{itemize}


% ============================================================
\section{当前文档索引}
% ============================================================

\subsection{核心产品文档（core/）}

\begin{table}[h]
\centering
\begin{tabular}{lll}
\toprule
\textbf{文件} & \textbf{标题} & \textbf{状态} \\
\midrule
MRD & 市场需求文档 & 已完成 \\
PRD & 产品需求文档 & 已完成 \\
RFC & 请求意见/技术提案 & 已完成 \\
TRD & 技术需求文档 & 已完成 \\
\bottomrule
\end{tabular}
\end{table}

\subsection{参考文档（reference/）}

\begin{table}[h]
\centering
\begin{tabular}{llll}
\toprule
\textbf{编号} & \textbf{文件名} & \textbf{标题} & \textbf{状态} \\
\midrule
REF-001 & skill-store-arch-survey & 开源 Skill 商店架构调研 & 已完成 \\
REF-002 & project-master-plan & 项目总纲（V0 先备阶段核心交付物） & 已完成 \\
REF-003 & open-skills-ecosystem & skills.sh 开放技能生态架构调研 & 已完成 \\
本文件 & doc-naming-convention & 文档命名规范 & 已完成 \\
\bottomrule
\end{tabular}
\end{table}

\subsection{需求文档（demands/）}

\begin{table}[h]
\centering
\begin{tabular}{llll}
\toprule
\textbf{文件} & \textbf{标题} & \textbf{类型} & \textbf{状态} \\
\midrule
MASTER & 总需求文档 & 总需求（持续更新） & v0.2 \\
DEMAND-000 & V0 先备阶段需求总纲 & 增量需求（归档） & 已完成 \\
DEMAND-001 & 角色分工重新定义 & 增量需求（归档） & 已完成 \\
DEMAND-002 & 开发环境搭建 & 增量需求（已完成） & 已完成 \\
DEMAND-003 & task-workflow & 任务流转路径定义 & 已完成 \\
DEMAND-004 & multi-agent-collaboration & 多智能体协作方案 & 已完成 \\
DEMAND-005 & mvp-upload-review & MVP 上传+审查可视化 & 已完成 \\
DEMAND-006 & frontend-rewrite & 前端重新开发（正式 UI） & 已完成 \\
DEMAND-007 & flow-fix-and-install & 功能流转修复与安装能力 & 已完成 \\
DEMAND-008 & dir-framework-rebuild & 目录框架重构与实习清单融合 & 已完成 \\
DEMAND-009 & email-auth-and-lifecycle & 邮箱认证注册+上架下架流程 & 进行中 \\
GUIDE-001 & project-structure & 项目目录结构与开发规范 & 已完成 \\
GUIDE-002 & project-methodology & 项目方法论总纲 & 已完成 \\
\bottomrule
\end{tabular}
\end{table}

\end{document}
